%!TEX root = ../main.tex
\section{Cross-Cutting Synthesis}
\label{sec:synthesis}

The operation taxonomy separates methods for analysis, but several conclusions only become visible when the stages are considered together. This section relates embodied preparation to language-model data engineering, resolves recurring conflicts across robot-data studies, and explains the field's movement from static curation toward feedback-driven preparation.

\subsection{What Transfers from Language-Model Data Preparation}

The most useful lesson from language-model data work is methodological. The pipeline should be treated as a versioned object that can be compared under controlled training. RefinedWeb and FineWeb disclose staged filtering and deduplication, while DataComp-LM fixes a candidate pool, model and compute tracks, and an evaluation suite so that pipeline choices can be tested as interventions~\cite{penedo2023refinedweb_2306_01116,penedo2024fineweb_2406_17557,li2024datacomp_2406_11794}. DoReMi and data selection for language-modeling via importance resampling (DSIR) optimize source weighting or resampling rather than hiding the mixture inside a training recipe~\cite{xie2023doremi_2305_10429,xie2023data_2302_03169}. LESS estimates how examples affect a target instruction-tuning objective, while Data-Efficient Instruction Tuning for Alignment (DEITA) separates quality, complexity, and diversity~\cite{xia2024less_2402_04333,liu2023what_2312_15685}. Joint-example selection (JEST) shows that selection can act on a jointly useful batch instead of ranking every example independently~\cite{evans2024data_2406_17711}.

Table~\ref{tab:llm-lessons} maps these mechanisms to embodied preparation. The mapping preserves the experimental idea while adding the properties that robot data require.

\begin{table*}[t]
\centering
\caption{Lessons from language and multimodal data engineering and the additional control required for embodied data.}
\label{tab:llm-lessons}
\scriptsize
\setlength{\tabcolsep}{4pt}
\begin{tabularx}{\textwidth}{@{}p{3.1cm}p{4.1cm}p{4.6cm}X@{}}
\toprule
Mechanism & Transferable lesson & Embodied adaptation & Additional control \\
\midrule
Fixed-pool competition~\cite{li2024datacomp_2406_11794,gadre2023datacomp_2304_14108}
& Hold the candidate pool, learner, compute, and evaluation fixed.
& Submit unit membership, weights, roles, lineage, and preparation cost.
& Separate simulation from blinded real-robot evaluation. \\
Auditable filtering~\cite{penedo2023refinedweb_2306_01116,penedo2024fineweb_2406_17557}
& Version each gate and deduplicate before splitting.
& Check technical validity, cross-modal alignment, and split leakage.
& Preserve temporal order and contact differences between visually similar episodes. \\
Mixture learning~\cite{xie2023doremi_2305_10429,xie2023data_2302_03169}
& Learn source weights against a declared target.
& Use hierarchical weights and protect rare but important groups.
& Test transfer from proxy learner or simulation to the target robot. \\
Targeted selection~\cite{xia2024less_2402_04333,liu2023what_2312_15685}
& Value depends on target, difficulty, and diversity.
& Select complementary units for a declared learner and task.
& Recalibrate after the policy changes the observed state distribution. \\
Joint batch selection~\cite{evans2024data_2406_17711}
& Example value depends on the surrounding batch.
& Select coherent skill or failure--recovery groups.
& Respect episode hierarchy and causal order. \\
Self-judged iteration~\cite{yuan2024self_2401_10020,shumailov2024collapse}
& Generation and judging can form a loop but may lose distribution tails.
& Retain audited real data while iterating on generated supervision.
& Use an independent verifier, mixture bounds, and rollback. \\
\bottomrule
\end{tabularx}
\end{table*}


A controlled robot-data competition should fix the raw pool, learner, training budget, and evaluation just as DataComp-LM controls its language setting. It must additionally preserve temporal and action semantics, separate simulation from physical execution, and account for reset, supervision, and safety costs. Robot transitions are not freely exchangeable tokens. Their meaning depends on controller, embodiment, timing, and the state reached by earlier actions.

Mixture learning also needs a hierarchical adaptation. Language corpora contain domains. Embodied corpora contain sources, tasks, environments, episodes, segments, and transitions. A source weight describes only the top of this hierarchy. A useful embodied method can learn coarse weights with a proxy learner, then refine routing at the task or segment level and test transfer to a larger VLA. Minimum sampling shares remain necessary for rare safety and recovery groups even when their short-term average loss is small.

Self-judged language-model iteration suggests a loop in which a model generates candidates, judges them, constructs training pairs, and updates itself~\cite{yuan2024self_2401_10020}. The corresponding embodied loop can generate alternatives, label rollout outcomes, and post-train a policy. Its risks are stronger because the output is an action. A model can approve an impossible contact, and a changing policy can shift the visual distribution on which its judge was calibrated. Work on recursively generated language data shows that indiscriminate replacement of natural data can erase parts of a distribution over generations~\cite{shumailov2024collapse}. This result does not prove robot-policy collapse. It motivates direct controls: retain audited real data, identify synthetic lineage, use independent checks, protect rare groups, and test retention after every round.

\subsection{Conditional Utility Resolves Conflicting Curation Results}

Formal analysis, learned mixture weights, subset-trained estimators, and influence methods all indicate that data value depends on a target learner or validation distribution~\cite{belkhale2023data_2306_02437,hejna2024re_2408_14037,dass2025datamil_2505_09603,lee2026quality_2603_09056}. This common dependence is more informative than whether a paper calls its score quality, progress, influence, or learnability. A selector without target evidence provides an initial estimate of usefulness, not a universal ranking.

The same view explains why strict filtering helps in some studies and hurts in others. Filtering can remove incorrect labels while also removing rare states, difficult tasks, and recoveries. Mutual-information filtering performs well in several settings but slightly underperforms the full set on a difficult Pen-in-Cup task~\cite{hejna2025robot_2502_08623}. Re-Mix reports a failure under an abnormal source, and a controlled curation audit shows that detecting an intended defect can disagree with policy utility~\cite{hejna2024re_2408_14037,bedi2026what_2606_10229}. The shared variable is the correctness--coverage balance left after selection. A method should report both dimensions and test the resulting policy under a fixed budget.

Mixed-quality data create a second apparent conflict. Removing failed demonstrations can improve imitation when every retained action is treated as expert behavior. Failure data can improve recovery or post-training when outcome, failure interval, retry, intervention, and desired behavior are represented. The $\pi_{0.7}$ metadata ablation, failure-segment methods, and retry-aware value learning support this distinction in different settings~\cite{intelligence20260_2604_15483,wu2024learning_2401_08957,qin2026beyond_2606_24633}. The disagreement is not whether a failure bit is good or bad. It is whether the preparation pipeline converts a mixed trajectory into units whose semantics match the training objective.

\subsection{The Useful Unit Is Often Smaller Than an Episode}

Segmentation, progress-based selection, retry modeling, and transition influence increasingly operate below the episode~\cite{wu2024learning_2401_08957,zhang2025scizor_2505_22626,wu2026sieve_2607_06442,qin2026beyond_2606_24633,lee2026quality_2603_09056}. This shift changes the question from whether to keep an episode to how to use its intervals. A valid approach, failed insertion, takeover, retry, and recovery can carry different targets and risks. Lineage must preserve their shared origin because segment value can depend on earlier actions and later outcome.

Smaller units do not imply independent scoring. An approach and grasp can be complementary, a failure and correction form a pair, and two individually useful sources can conflict in action convention. JEST's joint selection principle transfers here at the level of coherent behavior groups rather than arbitrary batches~\cite{evans2024data_2406_17711}. This suggests a hierarchy of decisions. Technical defects can be handled at frame or transition level, behavioral meaning at segment level, coverage at task or source level, and retention at the training-distribution level.

The same hierarchy clarifies generated-data claims. Observation editing, trajectory transformation, segment recomposition, simulation, and world-model generation change different parts of the supervision. Their validity conditions therefore differ. A plausible image may support representation learning, but action supervision needs consistent state transition, kinematics, contact, and target-controller semantics. MimicGen and Xiaomi-Robotics-U0 occupy different levels of this hierarchy and should not be compared by generated sample count alone~\cite{mandlekar2023mimicgen_2310_17596,li2026xiaomi_2607_11643}. The relevant statistic is how many units pass the checks required for their intended role and improve a policy under a matched training comparison.

\subsection{From Static Curation to Sequential Experimental Design}

The field's evolution is better explained by decision authority and feedback than by a list of publication dates. Static rules first made records technically usable. Pretrained models then scaled semantic labeling and screening. Learner-aware methods used target loss, influence, or subset training to estimate conditional value. Failure-driven post-training added states visited by the current policy. Data engines now repeat collection, verification, routing, training, and evaluation. The preparation agent in Figure~\ref{fig:autonomy-ladder} would also choose which experiment to run and when evidence is sufficient.

Each step increases target relevance but also couples the data more tightly to the current model. This creates a progression in both capability and risk. Static errors are repeatable. Learned filters introduce model bias. Learner-aware selection can overfit a target set. Closed loops can amplify the policy and verifier's shared blind spots. Autonomous operator choice can spend physical resources or remove coverage before its utility estimate is calibrated.

Our synthesis is that scalable preparation should be organized as sequential experimental design. Cheap deterministic checks remove known corruption. Model-based signals handle semantic cases at scale. Learner-aware estimates arbitrate when coverage and correctness conflict. Simulation and real execution calibrate the highest-impact decisions. A versioned controller chooses among these evidence sources under cost and safety limits. No single judge replaces the stack. The controller becomes more autonomous only as its decisions remain attributable, reversible, and supported by downstream evidence.
